(1) 单吸附位权重为空位的1及占据的 \(e^{-\beta(\varepsilon-\mu)}\)，故 \(\Xi_1=1+e^{\beta(\mu-\varepsilon)}\)，独立位点总巨配分函数 \(\Xi=\Xi_1^{N_s}\)。巨势 \(\Omega=-k_BT\ln\Xi\)，且 \(-\partial_\mu\Omega=\langle N\rangle\)。这里能级 \(\varepsilon\) 是单粒子吸附能，不应将 \(\mu\) 当作每个位都额外产生的系统内能。

(2) 占据概率 \(\vartheta=e^{\beta(\mu-\varepsilon)}/[1+e^{\beta(\mu-\varepsilon)}]=[1+e^{\beta(\varepsilon-\mu)}]^{-1}\)。因此 \(\langle N\rangle=N_s\vartheta\)，由于每位是参数为 \(\vartheta\) 的 Bernoulli 变量，\({\rm Var}(N)=N_s\vartheta(1-\vartheta)\)。其最大值在 \(\vartheta=1/2\) 时为 \(N_s/4\)，而相对涨落随系统大小按 \(N_s^{-1/2}\) 缩小。

(3) 对化学势求导，\(\partial\vartheta/\partial\mu=\beta e^{\beta(\varepsilon-\mu)}/[1+e^{\beta(\varepsilon-\mu)}]^2=\beta\vartheta(1-\vartheta)\)。所以 \(\partial\langle N\rangle/\partial\mu=\beta\,{\rm Var}(N)\)，是该巨正则系综的静态涨落--响应关系。曲线最陡处为 \(\mu=\varepsilon\)，其斜率 \(N_s/(4k_BT)\)，低温变陡。

(4) 独立二态的配置熵 \(S=-N_sk_B[\vartheta\ln\vartheta+(1-\vartheta)\ln(1-\vartheta)]\)，约定 \(0\ln0=0\)。当 \(\mu=\varepsilon\)，对所有 \(T>0\) 均有半覆盖 \(\vartheta=1/2\) 与最大熵 \(N_sk_B\ln2\)。固定 \(\mu\ne\varepsilon\) 且 \(T\to0^+\) 时，\(\vartheta\to1\) 若 \(\mu>\varepsilon\)，否则趋零。严格在 \(\mu=\varepsilon\) 的零温极限仍为半覆盖，这反映极限次序与基态简并。
